\chapter{Records and Immutability}
\label{ch:records}

\chapterorient{A simulation's state is rarely a single number. It is a
\emph{bundle}: the current field, how many iterations have run, whether the
solver has converged, the residual it is chasing. This chapter introduces the
data structure that carries such a bundle---the \emph{record}, a collection of
named, typed fields---and the one rule that governs every record in this book:
they never change. ``Updating'' a record does not overwrite it; it builds a
fresh one that differs in a field or two and leaves the original untouched. That
single discipline---update means fresh copy---is what turns a tangle of mutable
state into the clean, traceable dataflow the rest of Part~III is built on.}

\section{A bundle of named fields}

In \cref{ch:language} we built a vocabulary of \emph{values} and the
\emph{functions} that map between them. The values so far have been simple: a
tensor, an integer, a boolean. But the quantities a simulation actually carries
around are seldom so lonely. Consider what a linear solver knows about itself in
the middle of a run. It holds a current guess at the answer---a tensor. It has
taken some number of steps---an integer. It has decided, or not, that it is
close enough---a boolean. These three facts travel together; they are the state
of \emph{one} solver at \emph{one} moment. We want a single value that holds all
three.

That value is a \emph{record}: a bundle of named, typed fields. We write one
with braces, listing each field as \lstinline[style=l4style]|name: value| pairs
separated by commas. Here is the solver state just described:

\begin{lstlisting}[style=l4style]
{ x: theField, it: 7, converged: False }
\end{lstlisting}

This is a single value. It has three fields: \code{x}, whose value is some
tensor we have named \code{theField}; \code{it}, an integer that happens to be
$7$; and \code{converged}, a boolean that happens to be \code{False}. The names
\code{x}, \code{it}, \code{converged} are chosen by us and are part of the
record's \emph{type}, which we may write the same way but with types in the
value slots:

\begin{lstlisting}[style=l4style]
{ x: Tensor[N], it: Int, converged: Bool }
\end{lstlisting}

Read this as a sentence: ``a record with a field \code{x} holding a tensor of
length \code{N}, a field \code{it} holding an integer, and a field
\code{converged} holding a boolean.'' Every record value of this type has
exactly these three fields, with exactly these types. \Cref{fig:card} draws the
value as what it morally is: a small labeled card.

\begin{figure}[t]
\centering
\begin{tikzpicture}[font=\footnotesize]
  % the card
  \node[draw=simBlue, fill=simBgBlue, rounded corners=3pt, thick,
        minimum width=46mm, minimum height=30mm, inner sep=6pt] (card) at (0,0) {};
  \node[anchor=north, font=\scriptsize\bfseries, text=simBlue]
    at ($(card.north)+(0,-1.5mm)$) {SolverState};
  % rows
  \node[anchor=west] at ($(card.north west)+(3mm,-7mm)$)
    {\code{x}\,:\ \(\mathsf{Tensor}[N]\) \;\(=\ \mathbf{theField}\)};
  \node[anchor=west] at ($(card.north west)+(3mm,-14mm)$)
    {\code{it}\,:\ \(\mathsf{Int}\) \;\(=\ 7\)};
  \node[anchor=west] at ($(card.north west)+(3mm,-21mm)$)
    {\code{converged}\,:\ \(\mathsf{Bool}\) \;\(=\ \mathrm{False}\)};
  % separating rules
  \draw[simGray!40] ($(card.north west)+(2mm,-10mm)$) -- ($(card.north east)+(-2mm,-10mm)$);
  \draw[simGray!40] ($(card.north west)+(2mm,-17mm)$) -- ($(card.north east)+(-2mm,-17mm)$);
\end{tikzpicture}
\caption[A record as a labeled card]{A record is a card of named, typed fields.
Each row carries a field name, the field's type, and the value currently stored
there. The card as a whole is a single value of type
\code{\{ x: Tensor[N], it: Int, converged: Bool \}}; the rows
are reached by name.}
\label{fig:card}
\end{figure}

\subsection{Reaching a field: the dot}

To pull one field out of a record we write the record, a dot, and the field
name. If \code{s} is the record above, then

\begin{lstlisting}[style=l4style]
s.it          -- the integer 7
s.converged   -- the boolean False
s.x           -- the tensor theField
\end{lstlisting}

Field access \lstinline[style=l4style]|s.it| is a value-producing expression
like any other: it names the integer in the \code{it} slot. The reduction rule
is exactly what intuition demands---projecting a field from a record literal
returns the value in that slot:
\[
  \{l_1\!:\!v_1,\dots,l_n\!:\!v_n\}.\,l_k \;\to\; v_k .
\]
There is nothing more to it. A record is a finite map from names to values, and
the dot looks one up.

Field access composes with everything else in the language. Because
\lstinline[style=l4style]|s.it| is an ordinary expression, it can appear wherever
a value can: as a function argument \lstinline[style=l4style]|norm(s.x)|, in
arithmetic \lstinline[style=l4style]|s.it + 1|, in a comparison
\lstinline[style=l4style]|s.it < s.max_it|, or as the field of \emph{another}
record \lstinline[style=l4style]|{ count: s.it }|. A record is not a special kind
of thing you must unpack before use; its fields are values the moment you name
them, and they flow through the pure-function vocabulary of \cref{ch:language}
unchanged.

\begin{notationbox}
We write record \emph{types} and record \emph{values} with the same brace
syntax, distinguished by what sits after each colon. A \emph{type} has a type in
the value slot---\lstinline[style=l4style]|{ it: Int }|---while a \emph{value}
has an actual value---\lstinline[style=l4style]|{ it: 7 }|. Field access is the
dot: \lstinline[style=l4style]|s.it|. (When a field name collides with a TeX
underscore, as in \code{final\_res}, the underscore is escaped in running prose
but typed literally in code.)
\end{notationbox}

\subsection{Records versus tuples}

The calculus also has \emph{tuples}---ordered bundles written with parentheses,
\lstinline[style=l4style]|(theField, 7, False)|, whose components are reached by
\emph{position}: $\pi_1$ for the first, $\pi_2$ for the second, and so on. A
tuple and a record can carry the same data. The difference is how you address a
component, and the difference matters more than it first appears.

A tuple's components are anonymous. To read the iteration count from
\lstinline[style=l4style]|(theField, 7, False)| you must \emph{know} that the
count is the second slot---and if someone later inserts a field, every
$\pi_i$ downstream silently shifts to the wrong thing. A record's components
are named. \lstinline[style=l4style]|s.it| means the iteration count no matter
where in the record it sits, and adding a field disturbs nothing. For a bundle
that travels a long way through a program and accumulates fields over a project's
lifetime---which is exactly what a simulation's state does---names are
self-documenting and robust where positions are cryptic and brittle. We will use
tuples for short-lived, obviously-ordered pairs (a quotient and remainder, a
coordinate) and records for everything that earns a name.

\begin{keyidea}
A record is to a function what a \emph{struct} is to a procedure: a bundle of
named fields that a routine reads and ``updates.'' The whole of this chapter is
the one place that analogy breaks. A struct is a box you reach into and
overwrite; a record is a value you can only \emph{read}. To ``update'' a record
you build a new one. That single change---from a box you mutate to a value you
copy---is what the rest of Part~III rests on.
\end{keyidea}

\section{Immutable update: building a fresh copy}

Here is the central fact, and it is worth stating before any syntax. \emph{A
record, once built, never changes.} There is no operation that reaches into an
existing record and overwrites a field. None. The only thing you can do with a
record is read its fields and use them to build \emph{other} records. When we
speak---as we constantly will---of ``updating'' a record, we always mean: build
a new record, like the old one but differing in some field. The old record is
still there, unchanged, for anyone still holding it.

\subsection{The spread}

The syntax for ``a fresh record like \code{s} but with one field changed'' is the
\emph{spread}. We open a record literal, write \lstinline[style=l4style]|...s| to
mean ``copy in all of \code{s}'s fields,'' and then list the fields we want to
override:

\begin{lstlisting}[style=l4style]
{ ...s, it: s.it + 1 }
\end{lstlisting}

Read this aloud: ``a fresh record, with all of \code{s}'s fields, except that
\code{it} is \lstinline[style=l4style]|s.it + 1|.'' If \code{s} was
\lstinline[style=l4style]|{ x: theField, it: 7, converged: False }|, then the
expression above is a \emph{new} value
\lstinline[style=l4style]|{ x: theField, it: 8, converged: False }|. The fields
\code{x} and \code{converged} are copied straight across; \code{it} is replaced.
The reduction rule is precisely this---a spread over a record literal collapses
to a record literal with the named slot overwritten and the rest carried through:
\[
  \{\,\dots\{l_1\!:\!v_1,\dots,l_n\!:\!v_n\},\ l_k\!:\!w\,\}
  \;\to\;
  \{l_1\!:\!v_1,\dots,l_k\!:\!w,\dots,l_n\!:\!v_n\}.
\]

\Cref{fig:spread} draws what just happened. The original card \code{s} is on the
left, untouched. The spread produces a \emph{second} card on the right, identical
but for the one changed field. Both cards exist. Nothing was erased.

\begin{figure}[t]
\centering
\begin{tikzpicture}[font=\footnotesize]
  % --- left card: original s ---
  \begin{scope}
  \node[draw=simGray, fill=white, rounded corners=3pt, thick,
        minimum width=34mm, minimum height=26mm, inner sep=5pt] (s) at (0,0) {};
  \node[anchor=north, font=\scriptsize\bfseries, text=simGray] at ($(s.north)+(0,-1mm)$) {\code{s} (unchanged)};
  \node[anchor=west] at ($(s.north west)+(2.5mm,-7mm)$)  {\code{x}\,:\ \textbf{theField}};
  \node[anchor=west] at ($(s.north west)+(2.5mm,-13mm)$) {\code{it}\,:\ \(7\)};
  \node[anchor=west] at ($(s.north west)+(2.5mm,-19mm)$) {\code{converged}\,:\ F};
  \end{scope}
  % --- right card: fresh s' ---
  \begin{scope}[xshift=58mm]
  \node[draw=simBlue, fill=simBgBlue, rounded corners=3pt, thick,
        minimum width=34mm, minimum height=26mm, inner sep=5pt] (s2) at (0,0) {};
  \node[anchor=north, font=\scriptsize\bfseries, text=simBlue] at ($(s2.north)+(0,-1mm)$) {fresh copy};
  \node[anchor=west] at ($(s2.north west)+(2.5mm,-7mm)$)  {\code{x}\,:\ \textbf{theField}};
  \node[anchor=west, text=simRust, font=\footnotesize\bfseries]
        at ($(s2.north west)+(2.5mm,-13mm)$) {\code{it}\,:\ \(8\)};
  \node[anchor=west] at ($(s2.north west)+(2.5mm,-19mm)$) {\code{converged}\,:\ F};
  \end{scope}
  % the spread arrow
  \draw[flow] (s.east) -- node[above, font=\scriptsize\ttfamily]{\{...s, it: s.it+1\}}
                          node[below, font=\scriptsize, text=simGray]{copies \code{x}, \code{converged}} (s2.west);
\end{tikzpicture}
\caption[Spread produces a fresh card]{The spread
\code{\{ ...s, it: s.it + 1 \}} reads \code{s} and builds a
\emph{new} card. The unchanged fields \code{x} and \code{converged} are copied;
\code{it} is set to its successor (shown in rust). Crucially, the original card
\code{s} on the left is untouched---it is still a perfectly good value that
anyone holding it still sees as \code{it}~$=7$.}
\label{fig:spread}
\end{figure}

\subsection{Contrast: mutating a struct in place}

To feel the difference, picture the procedural alternative. In a language with
mutable structs you would write something like \code{s.it = s.it + 1}, and that
statement \emph{reaches into the box \code{s} and overwrites the \code{it} slot.}
After it runs, there is no ``old \code{s}'' anymore---the $7$ is gone, replaced
by $8$, in place. Anyone else holding a pointer to that same box now sees $8$
too, whether they expected the change or not. The assignment is a side effect: it
modifies the world, and its effect ripples to every name aliased to that box.

The spread does none of this. \lstinline[style=l4style]|{ ...s, it: s.it + 1 }|
is an \emph{expression}, not a command. It computes a value and changes nothing.
The $7$ in \code{s} is still $7$. The new value $8$ lives in a different
record, and the only way anyone sees it is by being handed that new record. This
is the difference between ``mutate the box'' and ``compute a new value,'' and it
is the line dividing the procedural world from the one we work in.

\begin{pitfall}
The single most common misreading of the spread is to think
\lstinline[style=l4style]|{ ...s, it: s.it + 1 }| \emph{changes \code{s}}. It
does not. The spread \emph{reads} \code{s} and \emph{produces a new record};
\code{s} is exactly as it was. If you find yourself writing ``and then \code{s}'s
\code{it} becomes $8$,'' stop---\code{s}'s \code{it} is still $7$ forever. What
became $8$ is the \code{it} of a different record, the one the spread built. The
moment you internalize that the spread never touches its source, the value-history
traces below read themselves.
\end{pitfall}

\subsection{Spreading several fields at once}

A spread may override more than one field; just list them all:

\begin{lstlisting}[style=l4style]
{ ...s, it: s.it + 1, converged: True }
\end{lstlisting}

This builds a fresh record advancing the counter \emph{and} flipping the
convergence flag in one stroke, copying \code{x} across unchanged. The order of
the overrides does not matter as long as no name is given twice; each names a
distinct slot of the result. And a field's new value may be computed from the old
record's fields---\lstinline[style=l4style]|s.it + 1| reads \code{s.it} to
compute the replacement---which is the normal way a step advances state: read the
current bundle, compute the next, package it as a fresh bundle.

\subsection{Reading a spread by its reduction rule}

Because the spread is just an expression with a reduction rule, we can evaluate
one by hand the way \cref{ch:language} evaluated function applications: rewrite
step by step until no rule applies. This is worth doing once explicitly---it
turns ``update means fresh copy'' from a slogan into a mechanical fact.

\begin{example}[A spread, reduced step by step]
\label{ex:reduce}
Take \lstinline[style=l4style]|s = { x: f, it: 7, converged: False }| and evaluate
the increment \lstinline[style=l4style]|{ ...s, it: s.it + 1 }|. First substitute
the value of \code{s} (the \textsf{let}-rule of \cref{ch:language}):
\begin{lstlisting}[style=l4style]
{ ...{ x: f, it: 7, converged: False }, it: s.it + 1 }
\end{lstlisting}
The field read \lstinline[style=l4style]|s.it| projects (the \textsf{proj} rule),
giving $7$, so the override value \lstinline[style=l4style]|s.it + 1| reduces to
$8$:
\begin{lstlisting}[style=l4style]
{ ...{ x: f, it: 7, converged: False }, it: 8 }
\end{lstlisting}
Now the spread rule fires: copy the inner record's fields, overwriting the named
one:
\begin{lstlisting}[style=l4style]
{ x: f, it: 8, converged: False }
\end{lstlisting}
No rule applies; this is the value. Notice what the trace makes undeniable: the
inner record \lstinline[style=l4style]|{ x: f, it: 7, converged: False }|---the
value of \code{s}---was \emph{read} (its fields copied, its \code{it} projected)
but at no point \emph{altered}. The result is a brand-new literal sitting beside
it. There is no rule in the calculus that overwrites a field of an existing
record; the spread rule \emph{constructs}, it does not assign.
\end{example}

This is why we can say, flatly, that records are immutable: not as a convention we
promise to honor, but because the only rules that act on records are projection,
destructuring, and spread, and \emph{none of them mutates}. Projection reads;
destructuring reads; spread builds anew. Immutability is a property of the
rules, not a discipline imposed on top of them.

\begin{workedexample}{Three updates, and the originals all survive}{history}
We start from an initial solver state and apply three spread-updates in sequence,
naming each result. The point is to watch every intermediate record \emph{survive}
its own update---a value-history trace.

\textbf{Setup.} Let
\begin{lstlisting}[style=l4style]
let s0 = { x: f0, it: 0, converged: False }
\end{lstlisting}
where \code{f0} is the initial field guess. Now three steps, each a spread:
\begin{lstlisting}[style=l4style]
let s1 = { ...s0, x: f1, it: 1 }
let s2 = { ...s1, x: f2, it: 2 }
let s3 = { ...s2, x: f3, it: 3, converged: True }
\end{lstlisting}
Here \code{f1}, \code{f2}, \code{f3} are the successively-improved fields (their
computation is the solver's business; \cref{ch:fold} supplies it). Each line
reads the previous record and builds a fresh one.

\tcblower
\textbf{The trace.} Because no spread ever mutates its source, all four records
coexist:
\[
\begin{array}{l@{\;=\;}l}
\code{s0} & \{\ \code{x}\!: f_0,\ \code{it}\!: 0,\ \code{converged}\!: \mathrm{F}\ \} \\[2pt]
\code{s1} & \{\ \code{x}\!: f_1,\ \code{it}\!: 1,\ \code{converged}\!: \mathrm{F}\ \} \\[2pt]
\code{s2} & \{\ \code{x}\!: f_2,\ \code{it}\!: 2,\ \code{converged}\!: \mathrm{F}\ \} \\[2pt]
\code{s3} & \{\ \code{x}\!: f_3,\ \code{it}\!: 3,\ \code{converged}\!: \mathrm{T}\ \}
\end{array}
\]
After all three updates, \lstinline[style=l4style]|s0.it| is still $0$ and
\lstinline[style=l4style]|s0.converged| is still \code{False}. The update that
produced \code{s1} did not reach back and alter \code{s0}; it produced a
\emph{new} record. Each later step likewise leaves its predecessor intact. We
have, for free, a complete history of the solve---every intermediate state, still
readable---simply because nothing was ever overwritten. \Cref{fig:chain} draws the
chain.
\end{workedexample}

\begin{figure}[t]
\centering
\begin{tikzpicture}[font=\footnotesize, node distance=12mm]
  \tikzset{rec/.style={draw, rounded corners=2pt, thick, align=center,
     minimum width=22mm, minimum height=13mm, inner sep=2pt, font=\scriptsize}}
  \node[rec, draw=simGray,  fill=white]     (a) {\code{s0}\\[-1pt]\(it\!=\!0\)\\\(\mathrm{F}\)};
  \node[rec, draw=simGray,  fill=white, right=of a] (b) {\code{s1}\\[-1pt]\(it\!=\!1\)\\\(\mathrm{F}\)};
  \node[rec, draw=simGray,  fill=white, right=of b] (c) {\code{s2}\\[-1pt]\(it\!=\!2\)\\\(\mathrm{F}\)};
  \node[rec, draw=simBlue, fill=simBgBlue, right=of c] (d) {\code{s3}\\[-1pt]\(it\!=\!3\)\\\(\mathrm{T}\)};
  \draw[flow] (a) -- node[above,font=\scriptsize]{spread} (b);
  \draw[flow] (b) -- node[above,font=\scriptsize]{spread} (c);
  \draw[flow] (c) -- node[above,font=\scriptsize]{spread} (d);
  \node[font=\scriptsize, text=simGray, align=center] at ($(a)+(0,-12mm)$)
    {\(it=0\)\\\emph{still readable}};
\end{tikzpicture}
\caption[A chain of fresh records]{The value history of
\cref{wex:history}. Each spread reads one record and emits a fresh successor;
the arrow is exactly the \emph{value-threading} of \cref{ch:targets}'s transient
solve. None of \code{s0}, \code{s1}, \code{s2} is destroyed when its successor is
built---the whole chain persists, so \code{s0.it} is $0$ even
after \code{s3} exists.}
\label{fig:chain}
\end{figure}

\section{Destructuring: pulling fields out by name}

Reaching fields one dot at a time is fine for a single field but clumsy when a
function needs several. \emph{Destructuring} lets a \code{let} bind several fields
to names in one move. Instead of binding the whole record to a name and dotting
into it, we write a brace pattern on the left of the \lstinline[style=l4style]|=|:

\begin{lstlisting}[style=l4style]
let { x, it } = s in
  -- here x is bound to s.x and it to s.it
  ...
\end{lstlisting}

This says: ``take \code{s}, and bind the local names \code{x} and \code{it} to its
\code{x} and \code{it} fields.'' Inside the body, \code{x} \emph{is}
\lstinline[style=l4style]|s.x| and \code{it} \emph{is} \lstinline[style=l4style]|s.it|,
without any further dotting. The reduction is the obvious one---destructuring a
record literal substitutes each bound name by the matching field:
\[
  \textsf{let}\ \{l_1,\dots,l_n\} = \{l_1\!:\!v_1,\dots,l_n\!:\!v_n\}\ \textsf{in}\ e
  \;\to\; e[v_1/l_1,\dots,v_n/l_n].
\]
A pattern need not mention every field---\lstinline[style=l4style]|let { it } = s|
binds only \code{it} and ignores the rest. \Cref{fig:destructure} pictures the
fields being pulled out by name.

\begin{figure}[t]
\centering
\begin{tikzpicture}[font=\footnotesize]
  \node[draw=simBlue, fill=simBgBlue, rounded corners=3pt, thick,
        minimum width=40mm, minimum height=26mm, inner sep=5pt] (s) at (0,0) {};
  \node[anchor=north, font=\scriptsize\bfseries, text=simBlue] at ($(s.north)+(0,-1mm)$) {\code{s}};
  \node[anchor=west] (fx)  at ($(s.north west)+(2.5mm,-7mm)$)  {\code{x}\,:\ theField};
  \node[anchor=west] (fit) at ($(s.north west)+(2.5mm,-13mm)$) {\code{it}\,:\ \(7\)};
  \node[anchor=west]       at ($(s.north west)+(2.5mm,-19mm)$) {\code{converged}\,:\ F};
  % pulled-out names
  \node[draw=simGreen, fill=simBgGreen, rounded corners=2pt, font=\ttfamily\footnotesize] (nx)  at (5.6,0.7) {x};
  \node[draw=simGreen, fill=simBgGreen, rounded corners=2pt, font=\ttfamily\footnotesize] (nit) at (5.6,-0.1) {it};
  \draw[flow, simGreen] (fx.east)  to[out=0,in=180] (nx.west);
  \draw[flow, simGreen] (fit.east) to[out=0,in=180] (nit.west);
  \node[font=\scriptsize, text=simGray, align=center] at (5.6,-1.0)
    {\ttfamily let \{ x, it \} = s};
\end{tikzpicture}
\caption[Destructuring pulls fields out by name]{The pattern
\code{let \{ x, it \} = s} binds the local names \code{x} and
\code{it} to the matching fields of \code{s}, by name. The \code{converged} field
is simply not mentioned, so no name is bound to it. Destructuring reads \code{s};
it does not consume or alter it.}
\label{fig:destructure}
\end{figure}

\subsection{Destructuring with a rest}

Sometimes a function wants a couple of fields by name \emph{and} the remaining
fields kept as a record---to forward them, or to spread them into a result. The
\emph{rest} pattern, written \lstinline[style=l4style]|...r|, captures whatever is
left:

\begin{lstlisting}[style=l4style]
let { it, ...rest } = s in
  { ...rest, it: it + 1 }
\end{lstlisting}

Here \code{it} is bound to \lstinline[style=l4style]|s.it|, and \code{rest} is
bound to a record holding all of \code{s}'s \emph{other} fields (\code{x} and
\code{converged}). The body then spreads \code{rest} back in and overrides
\code{it}---which is just our increment-step written so that the unchanged fields
flow through a name rather than being copied by \lstinline[style=l4style]|...s|.
The two idioms---\lstinline[style=l4style]|{ ...s, it: ... }| and destructure-with-rest
then re-spread---compute the same fresh record; pick whichever reads more clearly
at the use site.

\begin{workedexample}{Destructure, derive, rebuild---nothing mutated}{derive}
We destructure a state record, compute a derived quantity from its fields, and
build a \emph{new} record carrying that quantity---then check that the original
survived.

\textbf{Setup.} Suppose a step record carries a residual tensor and a target
tolerance:
\begin{lstlisting}[style=l4style]
let s = { r: residualVec, it: 4, tol: 1.0e-6 }
\end{lstlisting}
We want to decide convergence---is the residual's norm below \code{tol}?---and
record the verdict in a fresh state with the counter advanced.

\textbf{The step.} Destructure the fields we need, derive, rebuild:
\begin{lstlisting}[style=l4style]
let { r, it, tol } = s in
let nrm  = norm(r) in           -- derived scalar (a pure function of r)
let done = nrm < tol in         -- derived boolean
  { ...s, it: it + 1, converged: done }
\end{lstlisting}

\tcblower
\textbf{What happened.} The destructure read \code{r}, \code{it}, \code{tol} from
\code{s}. The two \code{let}s computed derived values---a norm and a comparison---
that are \emph{not stored in \code{s} at all}; they are ephemeral intermediates,
born here and gone after the result is built. The final spread built a fresh
record: all of \code{s}'s fields, with \code{it} advanced and a new
\code{converged} field set to \code{done}.

\textbf{The original survives.} After this expression evaluates,
\lstinline[style=l4style]|s| is byte-for-byte what it was:
\lstinline[style=l4style]|s.it| is still $4$, \lstinline[style=l4style]|s.r| is
still the same residual tensor, and \code{s} has no \code{converged} field at all.
Reading \code{s}, deriving from it, and building a new record from the derived
values \emph{never touched \code{s}}. This is the discipline in miniature: data
flows \emph{out} of records by destructuring and \emph{into} new records by
spreading, and the records in between are immutable all the way down.
\end{workedexample}

\section{Why immutability is the right discipline for a simulation}

The reader entitled to ask, at this point, why we tie our hands. Mutable structs
are the default in every systems language; copying a whole record to change one
field sounds wasteful. Three payoffs justify the discipline, and each one is a
load-bearing beam of the rest of the book.

\subsection{The dataflow is the dependency graph}

A simulation \emph{evolves} state: each step produces the next state from the
current one. When ``produces the next from the current'' literally means ``reads
record \code{s} and builds record \code{s'},'' the program's text \emph{is} a
graph of who-feeds-whom. The arrow from \code{s} to \code{s'} in
\cref{fig:chain} is not a metaphor; it is the dependency edge \code{s'} depends
on \code{s}, and it is visible right there in the spread that names \code{s}. In
\cref{ch:language} we saw that pure functions make data flow explicit; records
make the \emph{state} that flows be a first-class, inspectable value. Together
they make the whole computation a graph you can read off the page---which is
exactly what the calculus needs to reason about equivalence between layers, and
what a backend needs to schedule the work.

If instead each step \emph{mutated} a shared state box, the arrows would
vanish into invisible side effects. Two steps writing the same box would be
ordered only by accident of execution, and ``what does step~$k$ depend on'' would
be unanswerable from the text. Immutability buys us a readable graph; mutation
spends it.

\subsection{Only one field actually changes}

Look again at the increment step \lstinline[style=l4style]|{ ...s, it: s.it + 1 }|.
Of the three fields, \emph{one} is overridden; the other two are copied verbatim.
This is not an accident of the example. Across a great many solver steps, the
record that threads through carries several fields but the step touches
\emph{one}---the iteration counter ticks while the heavy field is recomputed
elsewhere, or the field advances while the counter and flags ride along. The
spread makes this visible: the override list is short, and it names exactly what
changed. We will return to this in \cref{ch:strata}, where the observation that a
solve's externally-visible state mutates at essentially a single point---the
counter---becomes the key to separating the strata of a record cleanly. For now,
notice only that the spread \emph{advertises} its single mutation point; a
mutable struct hides it among a dozen possible assignments.

\subsection{Sharing is safe}

Here is the payoff that the mutable world simply cannot offer. Because a record
never changes, two names may refer to the \emph{same} record with no hazard
whatsoever. There is no such thing as ``\code{a} aliases \code{b}, so changing
\code{a} surprises \code{b}''---because nothing can change \code{a}. Aliasing,
the perennial bug-farm of mutable state, is \emph{impossible to misuse}: the worst
that aliasing can do is let two names read the same unchanging value, which is
exactly what you want.

\begin{figure}[t]
\centering
\begin{tikzpicture}[font=\footnotesize]
  \node[draw=simBlue, fill=simBgBlue, rounded corners=3pt, thick, align=left,
        minimum width=36mm, minimum height=20mm, inner sep=4pt] (rec) at (0,0)
        {\scriptsize\code{x}: theField\\[1pt]\scriptsize\code{it}: \(7\)\\[1pt]\scriptsize\code{converged}: F};
  \node[font=\scriptsize\bfseries, text=simBlue, anchor=south] at (rec.north) {one immutable record};
  \node[draw=simGray, fill=white, rounded corners=2pt, font=\ttfamily\footnotesize] (na) at (-3.6,1.0) {a};
  \node[draw=simGray, fill=white, rounded corners=2pt, font=\ttfamily\footnotesize] (nb) at (-3.6,-1.0) {b};
  \draw[flow] (na.east) to[out=0,in=160] (rec.west);
  \draw[flow] (nb.east) to[out=0,in=200] (rec.west);
  \node[font=\scriptsize, text=simGreen, align=left, anchor=west] at (2.4,0)
    {both names read\\the same value;\\neither can\\change it};
\end{tikzpicture}
\caption[Safe sharing]{Two names \code{a} and \code{b} bound to the same
immutable record. In a mutable world this would be a hazard---a write through
\code{a} would silently change what \code{b} sees. Here it is free and safe:
because the record cannot change, sharing it is indistinguishable from each name
having its own copy, but costs nothing.}
\label{fig:sharing}
\end{figure}

\Cref{fig:sharing} draws the situation. This safety is not a minor convenience.
It is what lets a spread like \lstinline[style=l4style]|{ ...s, it: s.it + 1 }|
be cheap in practice: the new record need not deep-copy \code{s}'s tensor field
\code{x}: it can \emph{share} the very same immutable tensor, because that tensor
can never change underneath it. The conceptual model is ``fresh copy''; the
implementation, licensed by immutability, shares everything that did not change.
You get the clean semantics of copying with the runtime cost of pointer-sharing,
and you get it precisely \emph{because} nothing mutates.

\begin{keyidea}
Immutability turns aliasing from a hazard into a free optimization. Conceptually,
every spread builds a fresh record. In practice, the fresh record shares the
unchanged fields with its source---safely, because no one can ever mutate the
shared parts. ``Update means fresh copy'' is the \emph{semantics}; ``share what
didn't change'' is the \emph{implementation}, and immutability is exactly what
makes the second a sound realization of the first.
\end{keyidea}

\section{Records of records, optional fields, and demand}

Two refinements round out the record vocabulary; both will be used heavily later
and are introduced here only far enough to read.

\subsection{Records nest}

A field may itself hold a record. The simulator's full state is a record whose
fields are themselves bundles---a configuration sub-record, a solver sub-record,
an output sub-record. Access chains through dots,
\lstinline[style=l4style]|state.solver.it|, and a spread updates a nested field
by spreading at each level:

\begin{lstlisting}[style=l4style]
{ ...state, solver: { ...state.solver, it: state.solver.it + 1 } }
\end{lstlisting}

This builds a fresh outer record whose \code{solver} field is a fresh inner
record with \code{it} advanced---and every other field, inner and outer, copied
across. The nesting is just the rules we already have, applied at two levels.

There is a subtlety worth naming, because it is exactly where the procedural
intuition wants to lead you astray. To bump a nested counter you must rebuild
\emph{both} levels: the inner record (with \code{it} advanced) and the outer
record (with its \code{solver} field pointing at the new inner record). You
cannot ``reach in'' and bump \code{state.solver.it}---there is no such operation,
just as there is none at the top level. But this costs almost nothing in
practice: by the safe-sharing argument of the previous section, the
\emph{untouched} branches of \code{state}---its \code{cfg} field, the
\code{solver}'s other fields---are shared, not copied. Only the spine from the
root down to the changed field is rebuilt. A deep update touches one path through
the tree and shares everything off it, which is precisely the immutable-data
idiom every functional language uses to make ``copy the whole structure'' cheap.

\begin{example}[A nested update shares the untouched branch]
\label{ex:nested}
Let \lstinline[style=l4style]|state = { cfg: c, solver: { x: f, it: 2 } }|. The
expression \lstinline[style=l4style]|{ ...state, solver: { ...state.solver, it: 3 } }|
reduces to \lstinline[style=l4style]|{ cfg: c, solver: { x: f, it: 3 } }|. The
inner \code{x} field still holds the very same tensor \code{f} (shared, since it
did not change), and the outer \code{cfg} still holds the very same \code{c}.
Only two records were freshly built---the inner \code{solver} and the outer
\code{state}---and the original \code{state} is, as always, untouched and still
reports \lstinline[style=l4style]|state.solver.it| as $2$.
\end{example}

\subsection{Optional fields and demand-pruning}

A record type may declare a field \emph{optional}, written with a trailing
\lstinline[style=l4style]|?|:

\begin{lstlisting}[style=l4style]
{ x: Tensor[N], it: Int, residual?: Scalar }
\end{lstlisting}

The \code{residual?} field may or may not be present in a given value. Optionality
is the calculus's answer to a recurring need: a step can \emph{declare} an output
that downstream code might want---a residual norm for monitoring, say---without
forcing it to be computed when nobody asks. The rule, which \cref{ch:fold}
develops in full, is \emph{demand-pruning}: a field that no consumer reads is
never computed. \Cref{fig:optional} greys out exactly such a field.

\begin{figure}[t]
\centering
\begin{tikzpicture}[font=\footnotesize]
  \node[draw=simBlue, fill=simBgBlue, rounded corners=3pt, thick,
        minimum width=44mm, minimum height=26mm, inner sep=5pt] (s) at (0,0) {};
  \node[anchor=north, font=\scriptsize\bfseries, text=simBlue] at ($(s.north)+(0,-1mm)$) {step output};
  \node[anchor=west] at ($(s.north west)+(2.5mm,-7mm)$)  {\code{x}\,:\ theField};
  \node[anchor=west] at ($(s.north west)+(2.5mm,-13mm)$) {\code{it}\,:\ \(8\)};
  \node[anchor=west, text=simGray!60] at ($(s.north west)+(2.5mm,-19mm)$)
    {\itshape\code{residual?}\,:\ (not read \(\Rightarrow\) not computed)};
  \draw[simGray!40, densely dashed] ($(s.north west)+(2mm,-16mm)$) -- ($(s.north east)+(-2mm,-16mm)$);
\end{tikzpicture}
\caption[An optional, demand-pruned field]{A step declares an optional field
\code{residual?}. Because no downstream consumer reads it here, demand-pruning
elides it---the field is greyed out and never materializes. The same step,
consumed by code that \emph{does} read \code{residual?}, would compute it. The
declaration is fixed; whether the value is produced depends on demand.
\Cref{ch:fold} makes this precise.}
\label{fig:optional}
\end{figure}

This is a deep idea with a light footprint here: it replaces, in one stroke, the
logging flags, verbosity switches, and conditional ``if monitoring, also
compute\dots'' branches that clutter procedural simulators. An operator simply
declares every output it \emph{could} produce; consumers read what they need; the
unread outputs cost nothing. We meet it again, with its reduction rule, in
\cref{ch:fold}.

\section{The book's real records: a forward map}

Everything in this chapter is rehearsal for three records you will meet
everywhere in Part~III and beyond. We name them now---only their shapes, so that
when they appear they are familiar---and defer their full treatment to
\cref{ch:strata}, which exists precisely to dissect them.

\begin{definition}[The three workhorse records]
\label{def:workhorse}
The recurring state records of the simulator's solve are:
\begin{itemize}[nosep]
  \item \textbf{\code{OpParams}}---\emph{build-time configuration.} The readonly
  bundle a solver is set up with: variant selectors, tolerances, the constructed
  operator surfaces. Captured once, never re-evolved.
  \item \textbf{\code{SimState}}---\emph{run-time evolving state.} The
  caller-visible quantities a solve carries step to step: the current iterate, the
  iteration count, the convergence flag. The very record this chapter has been
  drawing as cards.
  \item \textbf{\code{SolveResult}}---\emph{a returned bundle.} The record one step
  hands back: the next state, the working data, and the demand-prunable per-step
  outputs.
\end{itemize}
\end{definition}

Their shapes, in the syntax of this chapter:

\begin{lstlisting}[style=l4style]
OpParams    = { tol: Float, max_it: Int, flexible: Bool, ... }   -- build-time, readonly
SimState    = { x: Tensor[N], it: Int, converged: Bool, ... }    -- run-time, evolving
SolveResult = { sim: SimState, outputs: StepOutputs, ... }       -- a step's returned bundle
\end{lstlisting}

Notice the strata at a glance. \code{OpParams} is fixed once and read forever---
the ``operator parameters'' a step closes over but never changes. \code{SimState}
is the bundle that threads, the one whose \code{it} ticks step by step.
\code{SolveResult} is what a step returns: it \emph{contains} a fresh
\code{SimState} and adds the per-step readouts. The discipline of this
chapter---read by destructuring, write by spreading, never mutate---is exactly how
every step over these records is written. \Cref{ch:strata} explains \emph{why} the
state splits into these three strata (build-time versus run-time, and the single
mutation point we flagged above); \cref{ch:monad} shows how the threading of
\code{SimState} through a chain of steps is itself packaged so the spreads stay
out of the algorithm's way. For this chapter it is enough that you can read each
shape as a card of named fields and know that ``updating'' any of them means
building a fresh one.

\summarysection
A \emph{record} is a bundle of named, typed fields,
written \lstinline[style=l4style]|{ x: theField, it: 7, converged: False }|, with
a field reached by the dot, \lstinline[style=l4style]|s.it|. Records carry the
state a simulation evolves. The governing discipline is \emph{immutability}: a
record never changes. ``Updating'' a record means building a fresh one with the
\emph{spread}, \lstinline[style=l4style]|{ ...s, it: s.it + 1 }|---a new value,
like \code{s} but for the named fields, with \code{s} left untouched. This
contrasts sharply with mutating a struct field in place, which destroys the old
value and ripples through every alias. Fields are pulled out by name with
\emph{destructuring}, \lstinline[style=l4style]|let { x, it } = s|, optionally
keeping a rest, \lstinline[style=l4style]|let { it, ...rest } = s|. Immutability
pays three dividends for simulation: the value-threading of fresh records
\emph{is} the dependency graph (\cref{ch:targets}); the short override list of a
spread advertises that typically a single field changes per step
(\cref{ch:strata}); and sharing a record is always safe, which lets ``fresh
copy'' be implemented as ``share what didn't change'' at no semantic cost. Records
nest, and optional fields (\lstinline[style=l4style]|residual?|) introduce
demand-pruning, by which an unread output is never computed (\cref{ch:fold}). The
book's three workhorse records---\code{OpParams} (build-time),
\code{SimState} (run-time), \code{SolveResult} (a step's return)---are previewed
here and dissected in \cref{ch:strata}.

\furtherreading
The record syntax used here---brace literals, field access by name, the spread
for non-destructive update, destructuring with rest---is the common surface of
the typed functional and modern scripting languages; the calculus borrows the
TypeScript spelling deliberately for readability. The deeper point, that
immutable values make sharing safe and dataflow explicit, is the thesis of pure
functional programming; the standard exposition, and the source of the
``everything is a value, updates build fresh values'' discipline this chapter
distills, is Peyton Jones's account of the lazy, purely functional
language~\citep{peytonjones2003}, where immutability and demand-driven evaluation
(our optional fields) appear as two faces of one design.

\exercisessection
\begin{enumerate}[nosep]
  \item Let \code{s = \{ x: a, it: 3, converged: False \}}.
  Write the value of each of \code{s.it},
  \code{\{ ...s, it: s.it + 1 \}}, and
  \code{\{ ...s, converged: True \}}. Then state the value of
  \code{s.it} \emph{after} all three expressions have been
  evaluated, and explain in one sentence why it is what it is.
  \item Rewrite the increment \code{\{ ...s, it: s.it + 1 \}}
  using destructure-with-rest instead of \code{...s}. Argue
  that the two produce the same record.
  \item A colleague writes, ``the spread \code{\{ ...s, it: 0 \}}
  resets \code{s}'s counter to zero.'' Correct the statement precisely. What, if
  anything, happens to \code{s}?
  \item Give a record value and a tuple value that carry the same three pieces of
  data. Then suppose a fourth field/component is inserted at the front. For each,
  say what (if anything) breaks in code that read the \emph{old} second
  field/component, and use this to explain when you would prefer a record over a
  tuple.
  \item Take \code{state = \{ cfg: c, solver: \{ x: f, it: 5 \} \}}.
  Write a single expression that produces a fresh \code{state} whose
  \code{solver.it} is $6$ and whose every other field (inner and outer) is
  unchanged.
  \item Explain why two names bound to the same record cause no aliasing hazard
  here, whereas two pointers to the same mutable struct do. Connect your answer
  to why a spread may safely \emph{share} the unchanged tensor field \code{x}
  rather than deep-copying it.
  \item A step record declares an optional field
  \code{residual?: Scalar}. Describe a consumer that causes it
  to be computed and a consumer that causes it to be pruned. In one sentence,
  state what procedural feature (flag, branch, or monad) the optional field
  replaces.
  \item In \cref{wex:history}, suppose a bug replaced the third step with an
  imagined in-place mutation \code{s2.it = 3; s2.converged = True} (not legal in
  our calculus). Which earlier records in the trace would now be wrong if other
  code still held them, and why? Use this to state, in your own words, the single
  guarantee immutability provides.
\end{enumerate}
